% RECOVERED_RESULT. II REV10, 05_accion.tex:1--195,
% revision_planck.tex:13--93 and definicion_planck_rev06.tex:4--34.
% Selection: action antecedents of the ellipse; electronic applications
% and the metrological comparison, which do not determine the radius, are excluded.
\section{Determinantal structure and oriented action sections}
\label{iv:action:section}

The readouts of \(\pi,\varphi,e,\alpha\) arise from the
APP--TRIT--TPK state of Sections
\ref{sec:nucleo}--\ref{sec:alpha}. The dodecaphasic record of
\ref{sec:k-registro} preserves its orientation, its remainders and
quotients, and the return memory
\eqref{eq:holonomia-memoria}. Two successive operations are defined on
this record: a local norm determines the valuation of the action
section, and a character with regional continuation determines its
coordinates before and after the return. Both operations transmit
their rules and normalizations to the angular construction.

\subsection{Local quotient and determinantal norm}
\label{iv:action:determinantal}

The local step-three block of the dodecaphasic record, identified
in \ref{subsec:k-hadamard}, is
\begin{equation}
 H_4=\begin{pmatrix}
 1&1&1&1\\
 1&1&-1&-1\\
 1&-1&1&-1\\
 1&-1&-1&1
 \end{pmatrix},
 \qquad G_H=\mathbb Z^4/H_4\mathbb Z^4.
 \label{iv:action:h4}
\end{equation}
The quotient corresponds to this local block. The three blocks of
the complete record retain their positions; the following norm
operation acts on the sheets of \(G_H\).

\begin{lemma}[Structure of the local quotient]
\label{iv:action:smith}
The Smith normal form of \(H_4\) is
\(\operatorname{diag}(1,2,2,4)\). Consequently,
\[
 G_H\simeq\mathbb Z/2\mathbb Z\oplus\mathbb Z/2\mathbb Z
       \oplus\mathbb Z/4\mathbb Z,\qquad |G_H|=16.
\]
\end{lemma}
\begin{proof}
The greatest common divisor of the entries is one. The minors of
order two are even and one has value \(-2\), so their greatest
common divisor is two. The identities
\(H_4H_4^{\mathsf T}=4I\) and \(\det H_4=-16\) show that the
cofactors equal \(\pm4\); the third determinantal divisor is
four. The fourth is \(16\). The successive quotients of
\(1,2,4,16\) are \(1,2,2,4\); their product determines the order
of the quotient.
\end{proof}

\begin{definition}[Determinantal action reader]
\label{iv:action:norma}
The scalar section \(\alpha\), constant on the sheets of \(G_H\),
and one application of the self-scaling factor \(\varphi\) determine
\begin{equation}
 \Sigma_H=\varphi\,\mathrm N_{G_H}(\alpha),\qquad
 \mathrm N_{G_H}(\alpha)=\prod_{g\in G_H}\alpha=\alpha^{16}.
 \label{iv:action:sigma}
\end{equation}
Its decimal valuation is
\[
 \operatorname{ord}_{10}(\Sigma_H)
   =\lfloor\log_{10}\Sigma_H\rfloor.
\]
\end{definition}
The sixteenth power follows from the number of sheets once the
multiplicative norm has been fixed. The norm and the single application
of \(\varphi\) are part of the stated reader: the group order,
considered in isolation, admits other functionals.

\begin{theorem}[Decimal valuation of the action section]
\label{iv:action:decada}
The preceding reader, applied to the HMT coordinates already generated,
satisfies
\begin{equation}
 10^{-34}<\Sigma_H<10^{-33},
 \qquad \operatorname{ord}_{10}(\Sigma_H)=-34.
 \label{iv:action:orden-decimal}
\end{equation}
\end{theorem}
\begin{proof}
We use the rational enclosures of the preceding readouts,
\begin{equation}
 \frac{1618}{1000}<\varphi<\frac{1619}{1000},
 \qquad
 \frac{7297}{10^6}<\alpha<\frac{7298}{10^6}.
 \label{iv:action:intervalos}
\end{equation}
The function \((x,a)\mapsto xa^{16}\) is increasing in both
positive variables. The integer inequalities
\[
 1618\cdot7297^{16}>10^{65},
 \qquad 1619\cdot7298^{16}<10^{66}
\]
give
\[
 10^{-34}
 <\frac{1618}{1000}\left(\frac{7297}{10^6}\right)^{16}
 <\Sigma_H
 <\frac{1619}{1000}\left(\frac{7298}{10^6}\right)^{16}
 <10^{-33}.
\]
The definition of the valuation completes the proof. The enclosures
are subsequent checks on the coordinates already produced; the
determinantal operation precedes any chart of units.
\end{proof}

\subsection{Action character and regional continuation}
\label{iv:action:caracter}

The central chart retains the matrix
\[
 B_c=\begin{pmatrix}7&2\\2&7\end{pmatrix},
 \qquad\lambda_+=9,\quad\lambda_-=5.
\]
The calendar has length \(12\), the step-three orbit has length
four, and the census quotient used is
\(104976/1944=54\), with the censuses defined in
Section~\ref{sec:extension}. These invariants determine the coefficients
assigned to degrees \(2,3,4,5\) of the character:
\begin{equation}
 \frac9{16}=\frac{\lambda_+}{4^2},\qquad
 \frac59=\frac{\lambda_-}{\lambda_+},\qquad
 \frac7{48}=\frac{\operatorname{tr}(B_c)/2}{4\cdot12},
 \qquad
 \frac1{54}=\frac{1944}{104976}.
 \label{iv:action:coeficientes}
\end{equation}
The assignment to these degrees and the alternating signs belong
to the definition of the analytic character. The operative content
combines the support, the orientation, and this assignment rule.

\begin{definition}[Local character and regional return reader]
\label{iv:action:lectores}
On the readouts \(\alpha,\varphi,\pi\), the fifth-order character,
the regional transfer, and its continuation are
\begin{align}
 H_5={}&\frac12\sqrt{\frac{1000\alpha}{\varphi}}-\alpha
       +\frac9{16}\alpha^2-\frac59\alpha^3
       +\frac7{48}\alpha^4-\frac1{54}\alpha^5,
       \label{iv:action:h5}\\
 D_A={}&\exp\!\left(-\frac{100\pi\alpha}{9}\right),
       \label{iv:action:da}\\
 \mathcal C_\pi(\alpha)={}&169\alpha^6+
       \frac{D_A\alpha^7}{1-D_A\alpha},
       \label{iv:action:cpi}\\
 \eta_{\mathrm{ret}}={}&H_5-\frac{90}{\pi}\mathcal C_\pi(\alpha).
       \label{iv:action:eta-ret}
\end{align}
The integer \(169\) is the regional selector of the readouts
constructed in \ref{sec:generacion}: the regional emission
\(169\,|\,272\,|\,272\) is retained, distinguished from
\(418\,|\,674\,|\,521\) in the exceptional construction.
The first strict prolongation of its continuation has weight
\(D_A\), and each concatenation multiplies the weight by
\(D_A\). The initial normalization is unity.
\end{definition}

The continuation is formulated before evaluation at \(z=\alpha\).
With \(D_A\) already produced by its reader, write
\(\mathscr C(z)=169z^6+\mathscr R(z)\), where
\(\mathscr R(z)\in z^7\mathbb R[[z]]\).
Concatenation increases the degree by one; the initial weight
and the transfer rule give
\begin{equation}
 \mathscr R(z)=D_Az^7+D_Az\mathscr R(z).
 \label{iv:action:renovacion}
\end{equation}

\begin{proposition}[Uniqueness of the normalized continuation]
\label{iv:action:unicidad-renovacion}
Equation \eqref{iv:action:renovacion} determines a unique series:
\begin{equation}
 \mathscr C(z)=169z^6+\frac{D_Az^7}{1-D_Az}
             =169z^6+\sum_{n=7}^{\infty}D_A^{n-6}z^n.
 \label{iv:action:serie-renovacion}
\end{equation}
Its realization is analytic for \(|D_Az|<1\).
\end{proposition}
\begin{proof}
The constant term of \(1-D_Az\) is one, so this element is
invertible in the ring of formal power series.
Solving \eqref{iv:action:renovacion} gives the rational expression.
Equivalently, the coefficient of degree seven is \(D_A\), and
each subsequent coefficient is \(D_A\) times its predecessor.
The geometric series converges absolutely in the stated disk.
\end{proof}

The normalization of the first weight enters the uniqueness statement.
The degree-six impulse and the concatenation ratio, without this
initial weight, are also compatible with
\[
 169z^6+\lambda\frac{D_Az^7}{1-D_Az}.
\]
The unit-normalization rule fixes \(\lambda=1\) before evaluation.
For \(0<\alpha<1\), we have \(0<D_A<1\) and \(D_A\alpha<1\);
therefore,
\begin{equation}
 \frac{D_A\alpha^7}{1-D_A\alpha}
   =\sum_{n=0}^{\infty}D_A^{n+1}\alpha^{n+7}.
 \label{iv:action:cola}
\end{equation}
After the term of index \(N\), the exact remainder is
\(D_A^{N+2}\alpha^{N+8}/(1-D_A\alpha)\). The rational expression
thus retains the complete continuation at any precision.

\subsection{Positive sections and action per revolution}
\label{iv:action:secciones}

Let \(\mathcal L_S\) be the oriented action line and
\(\mathcal U_S\) a positive basis. Its sections before and
after the return are
\begin{equation}
 \hbar_{\mathrm{pre}}=H_5\,10^{-34}\mathcal U_S,\qquad
 \hbar_{\mathrm{ret}}=\eta_{\mathrm{ret}}\,10^{-34}\mathcal U_S.
 \label{iv:action:hbar}
\end{equation}
The return preserves the regional compensation
\(90\mathcal C_\pi(\alpha)/\pi\) and the ratio
\begin{equation}
 \delta_{\mathrm{ret}}=H_5-\eta_{\mathrm{ret}},\qquad
 R_{\mathrm{act}}=\frac{\hbar_{\mathrm{pre}}}
                        {\hbar_{\mathrm{ret}}}
                =\frac{H_5}{\eta_{\mathrm{ret}}}.
 \label{iv:action:razon}
\end{equation}

\begin{proposition}[Positivity and covariance of the action ratio]
\label{iv:action:positividad}
On the intervals \eqref{iv:action:intervalos},
\(0<\eta_{\mathrm{ret}}<H_5\), and \(R_{\mathrm{act}}>1\).
The ratio remains fixed under a change of the common positive
basis \(\mathcal U_S\).
\end{proposition}
\begin{proof}
The terms of \(\mathcal C_\pi\) are positive. To bound
\(H_5\), it suffices to use \(\alpha>0.007\), \(\alpha<0.008\),
and \(\varphi<1.619\). The square-root term in \eqref{iv:action:h5}
is greater than one, and the sum of its negative terms is less
than \(0.008001\); hence \(H_5>0.99\).
With \(D_A<1\) and \(\pi>3\),
\[
 0<\frac{90}{\pi}\mathcal C_\pi
 <30\left(169(0.008)^6+
             \frac{(0.008)^7}{1-0.008}\right)<10^{-6}.
\]
Thus, \(\eta_{\mathrm{ret}}>0.989999>0\) and
\(\eta_{\mathrm{ret}}<H_5\). The change of basis divides both
action coordinates by the same positive factor, which cancels
in their ratio.
\end{proof}

The structural determination of angular action is the pair
\((\hbar_{\mathrm{pre}},\hbar_{\mathrm{ret}})\), with its
identification by return. It preserves the valuation of the local
norm, the action character, and the regional continuation.
In each section \(s\in\{\mathrm{pre},\mathrm{ret}\}\), the
action of a complete revolution is
\begin{equation}
 h_s=2\pi\hbar_s.
 \label{iv:action:vuelta}
\end{equation}
The subsequent chart \(\mathcal U_S\mapsto\mathrm{J\,s}\)
expresses these sections in units of action. The normalized radius
to be constructed remains independent of this common choice
of units.
